\section{Introduction}
\label{sec:intro}

\subsection{The problem}
\label{sec:intro-problem}

Let $G$ be a finite simple graph and let $A$ be an abelian group. An
\emph{$A$-flow} on $G$ is an assignment $\varphi$ of a value
$\varphi(e) \in A$ to every edge $e$, together with an orientation of
$e$, such that for every vertex $v$ the oriented sum
$\sum_{e \ni v}\sigma(v,e)\,\varphi(e)$ equals zero (Kirchhoff's
conservation law), where $\sigma(v,e) \in \{\pm 1\}$ records whether
$v$ is the head or tail of $e$. The flow is \emph{nowhere-zero} if
$\varphi(e) \ne 0$ for every $e$. The classical objects in this area
are nowhere-zero $k$-flows ($A = \Z$ and $\varphi(e) \in \{\pm 1, \dots,
\pm(k-1)\}$); a long-standing conjecture of Tutte asserts that every
bridgeless graph admits a nowhere-zero $5$-flow.

This paper is concerned with a continuous variant. An
\emph{$\Sphere$-flow} on $G$ assigns a unit vector $\varphi(e) \in \Sphere
\subset \R^{3}$ to every (oriented) edge so that Kirchhoff's law holds
at every vertex:
\begin{equation}
\label{eq:kirchhoff}
  \sum_{e \ni v} \sigma(v,e)\,\varphi(e) \;=\; 0 \qquad \text{for every } v \in V(G).
\end{equation}
The values are themselves unit vectors, so the nowhere-zero condition
is automatic. The natural existence question is:

\begin{quote}
  \textbf{Jain's first conjecture.} \emph{Every bridgeless graph
  admits an $\Sphere$-flow.}
\end{quote}

The conjecture is unsettled. For cubic graphs the Kirchhoff equation
at every vertex specialises to a constraint that three unit vectors
sum to zero; by an elementary calculation
(\cref{lem:120-degree}, after~\cite{HMM26}) this means the three are
coplanar with the origin and pairwise at angle $2\pi/3$.

\subsection{Jain's two conjectures and what is and is not known}
\label{sec:intro-disambiguation}

Jain proposed two conjectures, only the first of which is in scope
here. The second was a finite-labeling conjecture: there exists a map
$q : \Sphere \to \{\pm 1, \pm 2, \pm 3, \pm 4\}$ with $q(-x) = -q(x)$
and zero-sum on every great-circle triple of pairwise-equilateral
points. Were both conjectures true, they would jointly imply Tutte's
$5$-flow conjecture; this was the original motivation for the package.

Ulyanov~\cite{Uly26} recently \emph{disproved} the second conjecture
by exhibiting two finite point-set obstructions:
a 25-antipodal-pair instance whose nonexistence is certified by a SAT
instance with 200 variables and 19{,}765 clauses, and a 36-point
instance whose nonexistence at label range $\{\pm 1, \dots, \pm 5\}$
is certified by a 144-variable, 6710-clause SAT instance and verified
in Lean~4 via \texttt{bv\_decide}. The first conjecture is not touched
by Ulyanov's work; it remains the live problem. We will refer to it
simply as \emph{the conjecture} throughout.

A second relevant strand of recent work is the geometric
characterisation due to Houdrouge, Miraftab, and Morin~\cite{HMM26}.
They prove that for cubic graphs an $\Sphere$-flow is the same data as
an \emph{equiangular $\Sphere$-immersion} of the graph (every edge a
geodesic arc of fixed angular length, every face-flag of three edges
at a vertex meeting at angle $2\pi/3$). They also show that the
$\Sphere$-flow class is closed under \emph{triangle blow-up} of a
cubic vertex (the gadget we revisit and extend in
\cref{sec:gadget-closure}), and they give a rank obstruction: an
$\Sphere$-flow $\varphi$ with $\rank S_{\Q}(\varphi) \le 2$ and
odd-coordinate-free rational support induces a nowhere-zero $4$-flow.
The rank obstruction \emph{excludes} certain kinds of $\Sphere$-flow
on snarks but does not rule out arbitrary $\Sphere$-flows on them.

Closer to the route we follow, Mattiolo, Mazzuoccolo, Rajnik, and
Tabarelli~\cite{MMRT25} give an equivalence between oriented
cycle-double-covers and $H_d$-flows on cubic graphs, with the
consequence that an oriented $d$-CDC induces a nowhere-zero
$\Z_2^{d-2}$-flow and hence an $S^{d-2}$-flow. We will use the case
$d = 4$ as a sharp negative calibration in
\cref{sec:cdc-obstruction}.

\subsection{Our contributions}
\label{sec:intro-contributions}

We give three theorems, all with replayable computer-verified
certificates of the kind described in \cref{sec:certification}.

\begin{theorem}[Finite enumeration]
\label{thm:A-intro}
Every nontrivial snark on at most $28$ vertices admits an
$\Sphere$-flow. There are exactly $3\,247$ such graphs, with
per-order counts
\[
   |V| \;:\; 10\;(1),\; 18\;(2),\; 20\;(6),\; 22\;(20),\;
            24\;(38),\; 26\;(280),\; 28\;(2\,900),
\]
(no nontrivial snarks exist at any other order in this range). For
each we exhibit a numerical witness, refine it to $50$ decimal digits
in $\Q[s]/(s^{2}-3)$, and certify uniqueness in an explicit interval
box by the Krawczyk operator. The full certificate package is
replayable bit-for-bit; an independent cold-environment replay
re-verifies all $3\,247$ certificates.
\end{theorem}

\begin{theorem}[Constructive closure]
\label{thm:B-intro}
Let $\base$ be the set of $3\,247$ graphs of \cref{thm:A-intro} with
their certificates. Let $\closure$ be the closure of $\base$ under
\textbf{triangle blow-up} and \textbf{$3$-edge-cut dot product}. Every
$G \in \closure$ admits an $\Sphere$-flow that is produced
\emph{constructively}: the new flow is given in closed form from the
base witness(es) and the gadget data, and is verified directly by the
Kirchhoff and unit-norm equations at machine precision (residual at
most $5.9 \cdot 10^{-16}$ on every instance recorded here). No
per-graph Krawczyk solve is needed beyond the $3\,247$ base
certificates.
\end{theorem}

\begin{theorem}[Decomposition frontier]
\label{thm:C-intro}
Every graph in $\base$ is \emph{irreducible} under the two inverse
operations: there are no contractible triangles, and no cyclic
$3$-edge cuts (both verified by exhaustive scan). The catalogue
$\base$ is therefore exactly the base of $\closure$ under these
gadgets. Furthermore, for any cubic graph $G$, there is a memoised
dynamic program over canonical \texttt{graph6} strings
that returns a decomposition tree certifying
$G \in \closure$, or reports $G \notin \closure$ exhaustively.
\end{theorem}

We complement these positive results with three independent
\emph{negative} observations on the flower-snark family $J_{2k+1}$.
Each rules out one natural attempt to extend the gadget closure to a
uniform infinite family by structural means:
\begin{enumerate}[label=(\roman*)]
  \item a parity obstruction to a parametric (oriented) CDC on
        $J_{2k+1}$ for odd $n$ (\cref{sec:flower-cdc-fail});
  \item no $3$-edge-cut or $4$-edge-cut decomposition exists: $J_5$
        is cyclically exactly $5$-edge-connected and $J_{2k+1}$ for
        $k \ge 3$ is cyclically at least $6$-edge-connected (verified
        through $k = 6$), exceeding the textbook lower bound of $4$
        (\cref{sec:flower-no-dot-splice});
  \item the direct monodromy formulation has a transverse-mod-gauge
        fixed point on every certified $J_n$, but the LM-random
        witnesses do not converge to a smooth asymptotic limit and no
        stable small-period two-scale model fits them
        (\cref{sec:flower-monodromy-fail}).
\end{enumerate}

The flower question itself remains open; we do, however, conclude
that none of three standard structural attacks succeeds, and we leave
the diagnostic infrastructure available for future work.

\subsection{Why this package matters}
\label{sec:intro-why}

The finite theorem (\cref{thm:A-intro}) is, to our knowledge, the
largest interval-certified $\Sphere$-flow enumeration to date and a
natural baseline against which future symbolic or asymptotic
techniques can be compared. The closure theorem
(\cref{thm:B-intro}) is the first \emph{parametric} $\Sphere$-flow
construction that crosses outside the snark world while preserving
the geometric structure of~\cite{HMM26}: triangle blow-up shatters
girth $\ge 5$ and the dot product shatters cyclic $4$-edge-connectivity,
so $\closure$ properly extends $\base$ as a class of cubic graphs.
The decomposition theorem (\cref{thm:C-intro}) converts a positive
existence statement into a \emph{replayable} membership test: every
positive answer comes with a small operation tree that a third party
can re-execute against the base certificates.

The package is bookended by the prior work cited above: Jain's
formulation, Ulyanov's disambiguation, HMM's geometric characterisation,
and MMRT's CDC bridge. It does not settle Jain's first conjecture in
general; it does, however, give a publishable computational and
structural skeleton, with all the components either constructive or
ruled out by an explicit obstruction.

\subsection{Organisation}
\label{sec:intro-organisation}

\Cref{sec:preliminaries} fixes notation and recalls the geometric
characterisation. \Cref{sec:certification} describes the
interval-Krawczyk certifier and the schema-v2 replayable certificate
format. \Cref{sec:finite-theorem} states and proves
\cref{thm:A-intro}. \Cref{sec:cdc-obstruction} gives the CDC negative
calibration. \Cref{sec:gadget-closure} states and proves
\cref{thm:B-intro}. \Cref{sec:decomposition-frontier} states and
proves \cref{thm:C-intro}. \Cref{sec:flower-negatives} documents the
three independent obstructions on the flower family.
\Cref{sec:reproducibility} contains the reproducibility appendix.
